\section{Summary}\label{sec:synthesis}
A conic reformulation has several distinct costs. Local curvature bounds
the number and dimensions of the cone factors. Minimizing over whole
certificate fibers determines the least worst-case certificate rank;
globally $C^1$ selections of boundary points and certificates can raise
such minima. Affine restriction can lower the standard barrier
parameter; choosing another barrier on the same domain is a separate
optimization problem. For the cases in \eqref{eq:unsettled-range} and
for nonsymmetric barrier parameters, we prove only two-sided bounds. The integer tradeoffs are necessary
conditions whose minima need not be attained by lifts.

The weighted tree coefficient shows that certificate ranks do not
determine distance. In the squared leading coefficient, active nodes
contribute their full weight, and inactive nodes, whose subtrees miss
the support of the objective, contribute half, although their
certificate blocks vanish; along the central path, every node contributes
its full weight. The PSD Hessian theorem shows that exact allocation
equations improve the spectral comparison from squared to first powers
of the residual factors, but changing the auxiliary point can still
change the Newton right-hand side and the data it needs.

Geometric conclusions become computational ones only in a specified
model of access and output; a short description, a cheap normalized
state, a readable solution, and a reusable exact compiler are
different outputs. Under the compile-and-commit requirement of
Theorem~\ref{newt:dynamic-scale}, on an explicit family, the query cost
of maintaining scales grows linearly in the number of fresh update
batches; fixed data can instead be cached. Work bounds here state their access and per-round work assumptions, and metric bounds
their objective, endpoint, and barrier.
